\section{Proof of AI Work}

\label{sec:poai}
% ============================================================================

Proof of AI Work (PoAI) is a probabilistic verification protocol. Rather than
re-executing every inference request, verifiers sample a random subset and check
for consistency.

\subsection{Verification Protocol}

\begin{algorithm}
\caption{Proof of AI Work Verification}
\label{alg:poai}
\begin{algorithmic}[1]
\Require Inference records $\{\mathcal{I}_1, \ldots, \mathcal{I}_N\}$ from
    epoch $e$, sampling rate $p$
\Ensure Verification result for each sampled record
\State $S \leftarrow$ sample each $\mathcal{I}_i$ independently with
    probability $p$
\For{each $\mathcal{I}_i \in S$ in parallel}
    \State Retrieve input $\mathbf{x}_i$ from request log
    \State Re-execute: $\mathbf{y}_i' \leftarrow \texttt{Infer}(m_i,
        \mathbf{x}_i, \theta_i^{\text{det}})$
    \State Compare: $\delta_i = d(\mathbf{y}_i, \mathbf{y}_i')$
    \If{$\delta_i > \epsilon$}
        \State Flag worker $w_i$ for potential cheating
    \EndIf
\EndFor
\State Submit verification results to on-chain consensus
\end{algorithmic}
\end{algorithm}

\subsection{Handling Non-Determinism}

GPU floating-point arithmetic is non-deterministic across different hardware,
driver versions, and even consecutive runs on the same GPU. Two correct
executions of the same model on the same input may produce slightly different
outputs.

\begin{definition}[Verification Distance]
The verification distance $d(\mathbf{y}, \mathbf{y}')$ between two output
sequences uses a composite metric:
\begin{equation}
d(\mathbf{y}, \mathbf{y}') = \alpha \cdot d_{\text{token}}(\mathbf{y},
\mathbf{y}') + (1 - \alpha) \cdot d_{\text{logit}}(\mathbf{y}, \mathbf{y}'),
\end{equation}
where $d_{\text{token}}$ is normalized edit distance on generated tokens and
$d_{\text{logit}}$ is KL divergence on output logit distributions.
\end{definition}

To ensure deterministic verification, we fix the random seed and use
temperature $\tau = 0$ (greedy decoding) for verification re-execution:
\begin{equation}
\theta^{\text{det}} = (\tau = 0, \text{seed} = H(\mathcal{I}_i)),
\end{equation}
where the seed is derived from the inference record hash, making verification
reproducible.

\begin{theorem}[Detection Probability]
For a worker that cheats on fraction $\chi$ of requests, the probability of
detecting the cheating in a single epoch with sampling rate $p$ and $N$ total
requests is:
\begin{equation}
P_{\text{detect}} = 1 - (1 - p)^{\chi N}.
\end{equation}
For $p = 0.05$ (5\% sampling), $\chi = 0.10$ (cheating on 10\% of requests),
and $N = 10000$ (requests per epoch): $P_{\text{detect}} > 1 - 10^{-22}$.
\end{theorem}

\subsection{Economic Security}

\begin{definition}[Cheating Profitability]
A rational worker cheats if the expected profit from cheating exceeds the
expected loss from detection:
\begin{equation}
\chi \cdot R_{\text{save}} > P_{\text{detect}} \cdot S_{\text{slash}},
\end{equation}
where $R_{\text{save}}$ is compute cost saved per cheated request and
$S_{\text{slash}}$ is the slashing penalty.
\end{definition}

\begin{theorem}[Economic Security Bound]
Under the HANZO staking parameters~\cite{hanzotokenomics} (25\% slashing for
detected cheating, minimum stake proportional to GPU capacity), cheating is
unprofitable for any cheating rate $\chi > 0$ when the sampling rate satisfies:
\begin{equation}
p > \frac{\ln(R_{\text{save}} / S_{\text{slash}})}{\chi \cdot N}.
\end{equation}
With production parameters ($R_{\text{save}} / S_{\text{slash}} \approx 0.001$,
$N = 10000$): $p > 0.007$ suffices. We use $p = 0.05$ for a safety margin.
\end{theorem}

% ============================================================================
